% Start of the section on Biological Effects of Radiation
\section{EFECTOS BIOLÓGICOS DE LA RADIACIÓN}

% Subsection on Somatic and Genetic Effects
\subsection{Efectos Somáticos y Efectos Genéticos Radiosensibilidad}
Podrá señalar la diferencia entre los efectos somáticos y los efectos genéticos, así como la diferente radiosensibilidad de los distintos tipos de células, tejidos y órganos.

% Subsection on Dose-Response Relationship
\subsection{Relación Dosis-Efecto. Efecto de Dosis Aguda. Dosis Crónica y Efectos Tardíos.}
Describirá la diferencia entre dosis aguda y dosis crónica, distinguirá los valores de dosis que producen distintos efectos, desde un eritema hasta una dosis letal, así como la diferencia entre una irradiación localizada y una irradiación a cuerpo entero.

% Subsection on Forms of Exposure
\subsection{Formas de Exposición a la Radiación. Radiación Externa, Contaminación y Dosis Interna.}
El alumno distinguirá entre casos de exposición a radiación externa, interne, inhalación, y dosis interna.

% Start of the section on Radiation Medicine
\section{MEDICIÓN DE LA RADIACIÓN}

% Subsection on General Objective in Radiation Monitoring
\subsection{Objetivo General}
El alumno manejará diversos instrumentos de monitoreo, escogerá el adecuado para cada situación e interpretará correctamente las lecturas; quedará convencido de la ventaja de usar siempre un dosímetro personal, una alarma sonora y un instrumento de monitoreo, así como de cuidarlos adecuadamente.

% Subsection on Personal Monitoring Devices
\subsection{Dispositivos de Vigilancia Personal}
\subsubsection{Dosímetros de Bolsillo (0:30 hr)}
Entenderá el funcionamiento y manejará un dosímetro de bolsillo y su cargador, podrá verificar su funcionamiento y tomar las medidas pertinentes para su conservación. Interpretará correctamente las lecturas y conocerá sus ventajas y limitaciones.

% Practical Guide on the Use of Dosimeters
\subsection{PRÁCTICA NO. 2 - Manejo y uso del dosímetro de bolsillo de lectura directa (1:00 hr)}

\subsubsection{Dosímetros de Película y Termoluminiscentes (0:30 hr)}
El alumno reconocerá y podrá explicar la importancia del uso de dosímetros de registro temporal y la exposición acumulada. Se percatará de la necesidad de adquirir el hábito de usarlo y cuidar de que su uso y conservación sean adecuados.
